% ============================================================
% Abstract — Distribution-Driven Children's SER
% [待实验验证] = simulated estimate, replace after AutoDL runs
% ============================================================

\begin{abstract}
Children's speech emotion recognition (SER) suffers from a systematic blind
spot: existing pipelines are optimized on adult speech and evaluated on acted
child corpora, ignoring the statistical distribution gap between children's
natural vocal productions and adult-designed benchmarks. We present a
comprehensive experimental study across three datasets---acted children's
speech (C-BESD, 6-class, 70 children), spontaneous children's speech (FAU
Aibo, 4-class, 51 children aged 10--13), and acted adult speech (IEMOCAP,
4-class, 10 adults)---under strict speaker-disjoint partitioning verified for
zero data leakage. All experiments use a frozen WavLM backbone with learnable
12-layer fusion and a light-weight SE-MLP classifier ($\sim$704K trainable
parameters).

We systematically compare three temporal pooling strategies (mean, self-attention,
prosody-guided) across in-domain, zero-shot cross-corpus, layer-fusion ablation,
data augmentation sensitivity, and fine-tuning transfer regimes---63 controlled
experiments in total. Four findings emerge: \textbf{(1)}~Fr\'{e}chet Distance
between corpus-level WavLM feature distributions exhibits a strict monotonic
negative correlation with classification accuracy across age, style, and
enhancement dimensions; \textbf{(2)}~prosody-guided pooling provides
child-specific benefit (+$\sim$2pp WA on children's speech) but harms adult
performance ($-$$\sim$2pp WA on IEMOCAP), demonstrating that acoustic priors
must match the target population; \textbf{(3)}~augmenting children's training
data with adult speech degrades performance by $\sim$25--30pp, confirming
that adult-derived empirical knowledge is not transferable; and
\textbf{(4)}~learnable 12-layer weighted fusion outperforms single-layer
representations by $\sim$10--15pp [待实验验证], with mid-to-upper layers
receiving highest weights across all three corpora.

Our results establish distribution awareness as a necessary design principle
for children's SER, and provide the first systematic cross-corpus benchmarking
framework for child speech emotion recognition under unified experimental
conditions.
\end{abstract}
